\documentclass[11pt,a4paper]{article}

\usepackage{fontspec}
\IfFontExistsTF{Times New Roman}{\setmainfont{Times New Roman}}{\setmainfont{Tinos}}
\IfFontExistsTF{Consolas}{\setmonofont{Consolas}[Scale=0.88]}{\setmonofont{Liberation Mono}[Scale=0.88]}
\usepackage[english]{babel}
\usepackage{geometry}
\usepackage{xcolor}
\usepackage{enumitem}
\usepackage{listings}
\usepackage{booktabs}
\usepackage{array}
\usepackage{amsmath,amssymb}
\usepackage{titlesec}
\usepackage{fancyhdr}
\usepackage{parskip}
\usepackage{tcolorbox}
\tcbuselibrary{skins}
\usepackage{tabularx}
\usepackage{lastpage}
\usepackage{hyperref}

\geometry{top=1.35cm,bottom=1.35cm,left=1.65cm,right=1.65cm}

\definecolor{primaryblue}{RGB}{18,65,115}
\definecolor{secondaryteal}{RGB}{14,102,110}
\definecolor{correctgreen}{RGB}{10,120,40}
\definecolor{SoftBlue}{RGB}{232,242,250}
\definecolor{CodeBg}{RGB}{248,250,252}
\definecolor{DarkGray}{RGB}{55,55,55}
\definecolor{boxbg}{RGB}{248,250,252}

\setlength{\parindent}{0pt}
\setlength{\parskip}{0.8pt}
\setlist[itemize]{leftmargin=1.2em,itemsep=0.8pt,topsep=0.4pt,parsep=0pt}
\setlist[enumerate]{leftmargin=1.3em,itemsep=0.8pt,topsep=0.4pt,parsep=0pt}

\lstset{
  language=Python,
  basicstyle=\ttfamily\scriptsize,
  keywordstyle=\color{blue!75!black}\bfseries,
  stringstyle=\color{red!65!black},
  commentstyle=\color{primaryblue!80!black}\itshape,
  showstringspaces=false,
  columns=fullflexible,
  keepspaces=true,
  breaklines=true,
  frame=single,
  rulecolor=\color{primaryblue!45},
  backgroundcolor=\color{CodeBg},
  xleftmargin=0.2em,
  xrightmargin=0.2em,
  aboveskip=0.12em,
  belowskip=0.12em,
  tabsize=4
}

\pagestyle{fancy}
\fancyhf{}
\setlength{\headheight}{14pt}
\fancyhead[L]{\footnotesize\textbf{DSAI1003 -- Fundamental Programming Concepts in Python}}
\fancyhead[R]{\footnotesize\textbf{Solutions \& Marking Scheme (Loops \& Functions) --- Code: 202}}
\fancyfoot[L]{\scriptsize Faculty of Data Science and Artificial Intelligence -- National Economics University}
\fancyfoot[R]{\footnotesize Page \thepage\ of 6}
\renewcommand{\headrulewidth}{0.4pt}
\renewcommand{\footrulewidth}{0.3pt}

\titleformat{\section}{\large\bfseries\color{primaryblue}}{}{0pt}{}
\titleformat{\subsection}{\normalsize\bfseries\color{primaryblue}}{}{0pt}{}
\titlespacing*{\section}{0pt}{2.5pt}{1.0pt}
\titlespacing*{\subsection}{0pt}{1.8pt}{0.8pt}

\hypersetup{
  colorlinks=true,
  linkcolor=primaryblue,
  urlcolor=primaryblue,
  pdftitle={DSAI1003 - Midterm Exam Solutions Code 202 (Loops and Functions, 60 Minutes)},
  pdfauthor={Faculty of Data Science and Artificial Intelligence - National Economics University (NEU)}
}

\begin{document}

% ==============================================================================
% PAGE 1: HEADER + INSTRUCTOR NOTE + SECTION A PART 1 (Q1 - Q10)
% ==============================================================================
\begin{center}
  {\Large\textbf{NATIONAL ECONOMICS UNIVERSITY (NEU)}}\\[0.06cm]
  {\large\textbf{COLLEGE OF TECHNOLOGY}}\\[0.04cm]
  {\large\textbf{FACULTY OF DATA SCIENCE \& ARTIFICIAL INTELLIGENCE (FDA)}}\\[0.12cm]
  \rule{\textwidth}{1.2pt}\\[0.12cm]
  {\LARGE\bfseries\color{primaryblue}OFFICIAL SOLUTIONS \& MARKING RUBRIC}\\[0.06cm]
  {\large\textbf{Course: Fundamental Programming Concepts in Python (DSAI1003)}}\\[0.06cm]
  {\normalsize\textbf{Progress Midterm Exam: Loops, Iteration \& Modular Functions --- Code: 202}}\\[0.05cm]
  {\footnotesize\textbf{Weighting:} 40\% Quiz, 30\% Fill-in, 30\% Full Code \quad | \quad \textbf{Difficulty:} 80\% Basic (80 pts), 20\% Medium (20 pts)}\\[0.08cm]
  \rule{\textwidth}{0.6pt}
\end{center}

\vspace{0.05cm}

\begin{tcolorbox}[colback=green!8, colframe=correctgreen, arc=1.2mm, boxrule=0.6pt, left=2.0mm, right=2.0mm, top=0.8mm, bottom=0.8mm]
\small
\textbf{Note for Instructors and Teaching Assistants:}
This document contains official answer keys, conceptual justifications, reference implementations, sample outputs, and itemized grading rubrics for Code 202 of the 60-minute Loops \& Functions Midterm Exam. All grading must follow these rubrics consistently.
\end{tcolorbox}

\vspace{0.05cm}

\section*{Section A: Multiple-Choice Quiz Solutions (40 points --- 20 questions, 2 pts each)}
\textit{Difficulty: 16 Basic questions (32 pts, 80\%), 4 Medium questions (8 pts, 20\%).}

\subsection*{Part 1: Questions 1 to 10 (Loop Concepts \& Iteration Mechanics)}

\begin{center}
\renewcommand{\arraystretch}{1.10}
\setlength{\tabcolsep}{3.5pt}
\begin{tabularx}{\textwidth}{|c|c|c|X|c|}
\hline
\textbf{Q\#} & \textbf{Key} & \textbf{Level} & \textbf{Conceptual Explanation \& Rule Reference} & \textbf{Marks} \\
\hline
\textbf{Q1} & \textbf{B} & Basic & A \texttt{while} loop in Python is a pre-test loop: the boolean expression is evaluated before each iteration cycle, including prior to the first entry. & 2 pts \\
\hline
\textbf{Q2} & \textbf{B} & Basic & \texttt{range(10, 2, -2)} starts at 10, decrements by 2, and halts before reaching stop bound 2: values are $10, 8, 6, 4 \rightarrow$ \texttt{[10, 8, 6, 4]}. & 2 pts \\
\hline
\textbf{Q3} & \textbf{B} & Basic & \texttt{range(3, 8)} yields values $3, 4, 5, 6, 7$. Cumulative sum: $3 + 4 + 5 + 6 + 7 = 25$. & 2 pts \\
\hline
\textbf{Q4} & \textbf{B} & Basic & A sentinel is a distinctive, predefined input value used solely to indicate the end of data entry without being processed as a normal computation record. & 2 pts \\
\hline
\textbf{Q5} & \textbf{C} & Basic & In Python, \texttt{break} immediately terminates execution of the innermost enclosing loop and transfers control to the statement following the loop body. & 2 pts \\
\hline
\textbf{Q6} & \textbf{C} & Medium & Hand-trace odd values $m \in \{1, 3, 5, 7, 9, 11, 13, 15\}$: $m \% 3 == 0$ evaluates \texttt{True} for $m \in \{3, 9, 15\}$. Total sum: $3 + 9 + 15 = 27$. Loop terminates when $m = 17$. & 2 pts \\
\hline
\textbf{Q7} & \textbf{C} & Basic & Initializing \texttt{min\_val = 0.0} fails when all elements are positive. Initializing to the first item or positive infinity (\texttt{float('inf')}) guarantees correct comparisons. & 2 pts \\
\hline
\textbf{Q8} & \textbf{B} & Basic & In \texttt{"Py3.12-AI"}, \texttt{ch.isalpha()} evaluates \texttt{True} strictly for alphabetic characters: \texttt{'P'}, \texttt{'y'}, \texttt{'A'}, \texttt{'I'} (total: \texttt{4}). & 2 pts \\
\hline
\textbf{Q9} & \textbf{A} & Medium & Outer loop $a=1$: inner loop runs 1 time; $a=2$: 2 times; $a=3$: 3 times. Total iterations executed: $1 + 2 + 3 = 6$. & 2 pts \\
\hline
\textbf{Q10} & \textbf{B} & Basic & The \texttt{else} block of a loop executes if and only if the loop terminates normally (exhausting the iterable or condition becoming \texttt{False}), without encountering a \texttt{break}. & 2 pts \\
\hline
\end{tabularx}
\end{center}

\newpage

% ==============================================================================
% PAGE 2: SECTION A PART 2 + SECTION B INTRO & PROBLEM B1
% ==============================================================================
\subsection*{Part 2: Questions 11 to 20 (Functions, Modular Design \& Scope Mechanics)}

\begin{center}
\renewcommand{\arraystretch}{1.10}
\setlength{\tabcolsep}{3.5pt}
\begin{tabularx}{\textwidth}{|c|c|c|X|c|}
\hline
\textbf{Q\#} & \textbf{Key} & \textbf{Level} & \textbf{Conceptual Explanation \& Rule Reference} & \textbf{Marks} \\
\hline
\textbf{Q11} & \textbf{B} & Basic & Modular decomposition structures code into cohesive, reusable subroutines, eliminating code duplication and easing maintenance and unit testing. & 2 pts \\
\hline
\textbf{Q12} & \textbf{B} & Basic & \texttt{print()} writes human-readable output to the standard console stream, whereas \texttt{return} halts function execution and delivers a calculated value to the caller. & 2 pts \\
\hline
\textbf{Q13} & \textbf{A} & Basic & When a function is called, execution control temporarily suspends in the caller, transfers to the function block, and returns to the call site upon completion. & 2 pts \\
\hline
\textbf{Q14} & \textbf{C} & Basic & Parameters with default values must follow non-default parameters. In \texttt{def format\_record(name="Guest", dept, role="Analyst")}, non-default \texttt{dept} triggers a \texttt{SyntaxError}. & 2 pts \\
\hline
\textbf{Q15} & \textbf{B} & Medium & Invocation: \texttt{units = 6}, default \texttt{base = 12.0}, keyword \texttt{rate = 5.0}. Return value: $6 \times 5.0 + 12.0 = 30.0 + 12.0 = 42.0$. & 2 pts \\
\hline
\textbf{Q16} & \textbf{A} & Basic & Local variables have local scope; they exist only in the execution frame of the function and cannot be referenced from the global module scope. & 2 pts \\
\hline
\textbf{Q17} & \textbf{B} & Basic & When a local name matches a global name, Python resolves name lookups in the local frame first (LEGB rule), thereby shadowing the global identifier within the function. & 2 pts \\
\hline
\textbf{Q18} & \textbf{A} & Medium & Python integers are immutable. Inside \texttt{double(n)}, parameter \texttt{n} is rebound locally to $48$. Outer variable \texttt{num} remains unchanged at \texttt{24}. & 2 pts \\
\hline
\textbf{Q19} & \textbf{A} & Basic & Comma-separated expressions returned from a function (\texttt{return low, high}) are automatically packed into an immutable Python \texttt{tuple}. & 2 pts \\
\hline
\textbf{Q20} & \textbf{A} & Medium & High cohesion dictates that each function should encapsulate one single, clearly defined, and focused computational responsibility. & 2 pts \\
\hline
\end{tabularx}
\end{center}

\vspace{0.6mm}
\begin{tcolorbox}[colback=boxbg, colframe=primaryblue!45, arc=1.2mm, left=2.0mm, right=2.0mm, top=0.8mm, bottom=0.8mm]
\small
\textbf{Section A Distribution Summary:}\\
$\bullet$ Basic: 16 Qs $\times$ 2 pts (Q1--Q5, Q7--Q8, Q10--Q14, Q16--Q17, Q19) = \textbf{32 points (80.0\% of Section A)}\\
$\bullet$ Medium: 4 Qs $\times$ 2 pts (Q6, Q9, Q15, Q18, Q20) = \textbf{8 points (20.0\% of Section A)}\\
$\bullet$ Section A Total: \textbf{40 points (40.0\% of total exam score)}
\end{tcolorbox}

\vspace{-1.0mm}
\section*{Section B: Coding --- Fill-in Problems Solutions (30 points)}
\textit{Difficulty: 24 points Basic (80\%), 6 points Medium (20\%).}

\subsection*{Problem B1: Daily Temperature Logging \& Extreme Detection (12 points --- 4 blanks, 3 pts each [All Basic])}

\begin{lstlisting}
total_temp = 0.0
day_count = 0
temp = float(input("Enter temperature in Celsius (-999 to stop): "))  # [B1.1] (3 pts - Basic)
while temp != -999.0:                                               # [B1.2] (3 pts - Basic)
    total_temp = total_temp + temp                                  # [B1.3] (3 pts - Basic)
    day_count = day_count + 1
    temp = float(input("Enter next temperature (-999 to stop): "))
print(f"Total: {total_temp:.1f} C recorded across {day_count} days") # [B1.4] (3 pts - Basic)
\end{lstlisting}

\begin{center}
\renewcommand{\arraystretch}{1.12}
\setlength{\tabcolsep}{3.5pt}
\begin{tabularx}{\textwidth}{|c|c|X|c|}
\hline
\textbf{Blank} & \textbf{Level} & \textbf{Accepted Answer(s) \& Grading Criteria} & \textbf{Marks} \\
\hline
\textbf{[B1.1]} & Basic & \texttt{float(input("Enter temperature in Celsius (-999 to stop): "))} or \texttt{float(input())}\newline
\textit{Award 3 pts for correct float cast and input call. Award 1.5 pts if \texttt{int} cast used.} & 3 pts \\
\hline
\textbf{[B1.2]} & Basic & \texttt{temp != -999.0} or \texttt{temp != -999} or \texttt{temp > -999.0}\newline
\textit{Award 3 pts for valid sentinel inequality condition.} & 3 pts \\
\hline
\textbf{[B1.3]} & Basic & \texttt{total\_temp + temp}\newline
\textit{Award 3 pts. Accept \texttt{total\_temp += temp} if written as an assignment statement.} & 3 pts \\
\hline
\textbf{[B1.4]} & Basic & \texttt{total\_temp} (since \texttt{:.1f} follows in template) or \texttt{total\_temp:.1f}\newline
\textit{Award 3 pts for providing the correct accumulator variable name.} & 3 pts \\
\hline
\end{tabularx}
\end{center}

\newpage

% ==============================================================================
% PAGE 3: SECTION B PROBLEM B2 + SECTION B SUMMARY
% ==============================================================================
\subsection*{Problem B2: Modular Employee Sales \& Commission Processor (18 points --- 6 blanks, 3 pts each)}

\begin{lstlisting}
def calc_commission(sales, rate=0.08):                              # [B2.1] (3 pts - Basic)
    return sales * rate

sales_records = [1200.0, 4500.0, 8000.0, 2300.0, 6000.0]
total_commission = 0.0
max_sales = 0.0

for s in sales_records:                                             # [B2.2] (3 pts - Basic)
    comm = calc_commission(s)                                      # [B2.3] (3 pts - Basic)
    total_commission = total_commission + comm
    if s > max_sales:
        max_sales = s                                               # [B2.4] (3 pts - Basic)

bonus_pool = calc_commission(total_commission, rate=0.15)           # [B2.5] (3 pts - Medium)
avg_commission = total_commission / len(sales_records) if len(sales_records) > 0 else 0.0  # [B2.6] (3 pts - Medium)
print(f"Bonus Pool: ${bonus_pool:.2f} | Max Sales: ${max_sales:.2f}")
\end{lstlisting}

\begin{center}
\renewcommand{\arraystretch}{1.15}
\setlength{\tabcolsep}{3.5pt}
\begin{tabularx}{\textwidth}{|c|c|X|c|}
\hline
\textbf{Blank} & \textbf{Level} & \textbf{Accepted Answer(s) \& Grading Criteria} & \textbf{Marks} \\
\hline
\textbf{[B2.1]} & Basic & \texttt{rate=0.08}\newline
\textit{Award 3 pts for declaring default parameter with rate 0.08.} & 3 pts \\
\hline
\textbf{[B2.2]} & Basic & \texttt{sales\_records}\newline
\textit{Award 3 pts for iterating directly through the list collection.} & 3 pts \\
\hline
\textbf{[B2.3]} & Basic & \texttt{calc\_commission(s)}\newline
\textit{Award 3 pts for calling helper function with loop variable \texttt{s}.} & 3 pts \\
\hline
\textbf{[B2.4]} & Basic & \texttt{s}\newline
\textit{Award 3 pts for updating \texttt{max\_sales} with current sale amount.} & 3 pts \\
\hline
\textbf{[B2.5]} & Medium & \texttt{calc\_commission(total\_commission, rate=0.15)} or \texttt{calc\_commission(total\_commission, 0.15)}\newline
\textit{Award 3 pts for re-invoking helper function with bonus rate argument.} & 3 pts \\
\hline
\textbf{[B2.6]} & Medium & \texttt{total\_commission / len(sales\_records) if len(sales\_records) > 0 else 0.0}\newline
\textit{Award 3 pts for safe average calculation. Accept \texttt{total\_commission / len(sales\_records)} with 1.5 pts.} & 3 pts \\
\hline
\end{tabularx}
\end{center}

\vspace{1mm}
\begin{tcolorbox}[colback=boxbg, colframe=primaryblue!45, arc=1.2mm, left=2.0mm, right=2.0mm, top=1.0mm, bottom=1.0mm]
\small
\textbf{Section B Summary:}\\
$\bullet$ Basic: [B1.1]--[B1.4] (12 pts) + [B2.1]--[B2.4] (12 pts) = \textbf{24 points (80.0\% of Section B)}\\
$\bullet$ Medium: [B2.5] (3 pts) + [B2.6] (3 pts) = \textbf{6 points (20.0\% of Section B)}\\
$\bullet$ Section B Total: \textbf{30 points (30.0\% of total exam)}
\end{tcolorbox}

\newpage

% ==============================================================================
% PAGE 4: SECTION C PROBLEM C1 & PROBLEM C2
% ==============================================================================
\section*{Section C: Applied Full-Code Programming \& Design Solutions (30 points)}
\textit{Difficulty Distribution: 24 points Basic (80\%), 6 points Medium (20\%). Problems C1 and C2 share the common logistics fleet scenario. Problem C3 evaluates a sentinel loop and transaction auditing.}

\begin{tcolorbox}[colback=boxbg, colframe=secondaryteal, arc=1.2mm, boxrule=0.6pt, left=2.0mm, right=2.0mm, top=0.6mm, bottom=0.6mm]
\textbf{\small Problem Scenario for C1 \& C2: Logistics Delivery Fleet Fuel \& Trip Cost Analyzer}\\[0.1mm]
\footnotesize
A delivery logistics company tracks vehicle performance across $N$ delivery trips based on the following rules:
\begin{itemize}[itemsep=0.1pt,topsep=0.1pt]
  \item \textbf{Formulas per Trip:}
    $\text{Fuel Used (L)} = \text{distance\_km} \times (\text{liters\_per\_100km} / 100.0)$; \quad
    $\text{Trip Cost (USD)} = \text{Fuel Used} \times 1.45$.
  \item Across all $N$ trips, tally total distance (km), total fuel used (L), total cost (USD), and maximum single-trip distance.
\end{itemize}
\end{tcolorbox}

\vspace{0.2mm}
\subsection*{Problem C1: Algorithmic Specification \& Pseudocode (8 points --- [All Basic])}

\begin{lstlisting}[language={}]
BEGIN Fleet_Cost_Analyzer
    INPUT total_trips
    SET total_dist = 0.0, total_fuel = 0.0, total_cost = 0.0, max_dist = 0.0
    FOR trip FROM 1 TO total_trips DO:
        INPUT distance_km, liters_per_100km
        SET fuel_used = distance_km * (liters_per_100km / 100.0)
        SET trip_cost = fuel_used * 1.45
        SET total_dist = total_dist + distance_km
        SET total_fuel = total_fuel + fuel_used
        SET total_cost = total_cost + trip_cost
        IF distance_km > max_dist THEN:
            SET max_dist = distance_km
        END IF
    END FOR
    PRINT total_dist, total_fuel, total_cost, max_dist
END
\end{lstlisting}

\vspace{0.2mm}
\subsection*{Problem C2: Modular Python Function Implementation (10 points --- [All Basic])}

\begin{lstlisting}
def calculate_trip_stats(distance_km, liters_per_100km, price_per_liter=1.45):
    """Computes fuel used (liters) and cost (USD) for a single vehicle trip."""
    fuel_used = distance_km * (liters_per_100km / 100.0)
    cost = fuel_used * price_per_liter
    return fuel_used, cost

def main():
    n = int(input("Enter number of delivery trips: "))
    total_distance = 0.0
    total_fuel = 0.0
    total_cost = 0.0
    
    for i in range(1, n + 1):
        dist = float(input(f"Trip {i} distance (km): "))
        rate = float(input(f"Trip {i} fuel rate (L/100km): "))
        fuel, cost = calculate_trip_stats(dist, rate)
        print(f"Trip {i}: Distance = {dist:.2f} km | Fuel = {fuel:.2f} L | Cost = ${cost:.2f}")
        total_distance += dist
        total_fuel += fuel
        total_cost += cost
        
    print("-" * 56)
    print(f"Total Trips: {n} | Total Distance: {total_distance:.2f} km | Total Cost: ${total_cost:.2f}")

if __name__ == "__main__":
    main()
\end{lstlisting}

\begin{center}
\renewcommand{\arraystretch}{1.06}
\setlength{\tabcolsep}{3.5pt}
\begin{tabularx}{\textwidth}{|X|c|c|}
\hline
\textbf{Combined Grading Rubric for Problems C1 \& C2 (18 points total)} & \textbf{Level} & \textbf{Marks} \\
\hline
\textbf{[C1] Algorithmic Flow:} Proper input initialization and iterative loop header traversing all $N$ trips & Basic & 4.0 pts \\
\hline
\textbf{[C1] Computation Logic:} Correct formula evaluation, running totals accumulation, and simple max comparison & Basic & 4.0 pts \\
\hline
\hline
\textbf{[C2] Function Design:} Correct signature, parameters, arithmetic formulas, and tuple return & Basic & 4.0 pts \\
\hline
\textbf{[C2] Driver Loop \& IO:} Prompts for $N$, loops $N$ times, invokes function, unpacks tuple, accumulates, formats output & Basic & 6.0 pts \\
\hline
\end{tabularx}
\end{center}

\newpage

% ==============================================================================
% PAGE 5: PROBLEM C3 IMPLEMENTATION & VERIFICATION TRACE
% ==============================================================================
\subsection*{Problem C3: Retail Store Daily Transaction Auditor \& High-Value Sale Analyzer (12 points)}
\textit{Difficulty Breakdown: 6 points Basic + 6 points Medium.}

\textbf{Reference Python Implementation:}
\begin{lstlisting}
def is_high_value(amount, threshold=100.0):
    """Returns True if transaction amount meets or exceeds threshold, False otherwise."""
    return amount >= threshold

def main():
    SENTINEL = -1.0
    count = 0
    high_count = 0
    total_revenue = 0.0
    highest_sale = float("-inf")
    
    # Priming read
    amount = float(input("Enter transaction amount (-1.0 to finish): "))
    
    # Sentinel-controlled loop
    while amount != SENTINEL:
        count += 1
        total_revenue += amount
        
        if amount > highest_sale:
            highest_sale = amount
            
        if is_high_value(amount):
            high_count += 1
            
        amount = float(input("Enter next amount (-1.0 to finish): "))
        
    # Summary reporting
    if count == 0:
        print("No transactions were recorded.")
        return
        
    avg_amount = total_revenue / count
    high_rate = (high_count / count) * 100.0
    
    print("\n" + "=" * 56)
    print("DAILY TRANSACTION AUDIT SUMMARY")
    print("=" * 56)
    print(f"Total Transactions: {count} | High-Value: {high_count} ({high_rate:.1f}%)")
    print(f"Average Amount:     ${avg_amount:.2f} | Highest Sale: ${highest_sale:.2f}")
    print("=" * 56)

if __name__ == "__main__":
    main()
\end{lstlisting}

\vspace{0.4mm}
\textbf{Test Execution Trace \& Sample Output Verification:}
\begin{lstlisting}[basicstyle=\ttfamily\scriptsize]
Input Sequence: 45.0, 150.0, 85.0, 350.0, 82.5, -1.0
Execution Trace:
- Valid transaction count = 5
- Total revenue = 45.0 + 150.0 + 85.0 + 350.0 + 82.5 = 712.50 -> Average = $142.50
- High-value sales (>= $100.0) = 3 (150.0, 350.0, 82.5 is False; counts = 150, 350) -> Rate = 60.0%
- Highest single sale = $350.00
\end{lstlisting}

\newpage

% ==============================================================================
% PAGE 6: PROBLEM C3 RUBRIC & GRAND TOTAL SUMMARY TABLE
% ==============================================================================
\subsection*{Itemized Grading Rubric for Problem C3 (12 points)}

\begin{center}
\renewcommand{\arraystretch}{1.18}
\setlength{\tabcolsep}{3.5pt}
\begin{tabularx}{\textwidth}{|X|c|c|}
\hline
\textbf{Assessment Dimension \& Technical Criteria} & \textbf{Level} & \textbf{Marks} \\
\hline
\textbf{Modular Helper Function:} Defines \texttt{is\_high\_value(amount, threshold=100.0)} with default parameter and returns boolean outcome \texttt{amount >= threshold}. & Basic & 3.0 pts \\
\hline
\textbf{Sentinel Loop Mechanics:} Implements priming read, correct \texttt{while} loop termination on \texttt{-1.0}, and updates next read at the end of each iteration. & Basic & 3.0 pts \\
\hline
\textbf{Statistical Aggregators:} Accurately tallies total transaction count, running total revenue, high-value count (calling \texttt{is\_high\_value}), and dynamically tracks highest sale. & Medium & 3.0 pts \\
\hline
\textbf{Zero-Data Guard \& Formatted Reporting:} Includes zero-transaction guard (\texttt{count == 0}), computes average and high-value percentage, and prints formatted summary report. & Medium & 3.0 pts \\
\hline
\end{tabularx}
\end{center}

\vspace{2mm}
\begin{tcolorbox}[colback=primaryblue!6, colframe=primaryblue, arc=1.5mm, boxrule=0.8pt, left=3.0mm, right=3.0mm, top=1.8mm, bottom=1.8mm]
\small
\textbf{\large Overall Examination Structure, Weighting \& Difficulty Verification}\\[1.5mm]
\renewcommand{\arraystretch}{1.15}
\begin{tabularx}{\linewidth}{|X|c|c|c||c|}
\hline
\textbf{Examination Component} & \textbf{Basic (80\%)} & \textbf{Medium (20\%)} & \textbf{Total Points} & \textbf{Exam Weight} \\
\hline
\textbf{Section A: Multiple-Choice Quiz} & 32 pts (16 Qs) & 8 pts (4 Qs) & 40 pts & 40\% \\
\hline
\textbf{Section B: Fill-in Problems (B1 \& B2)} & 24 pts (8 blanks) & 6 pts (2 blanks) & 30 pts & 30\% \\
\hline
\textbf{Section C: Programming \& Design (C1--C3)} & 24 pts & 6 pts & 30 pts & 30\% \\
\hline
\hline
\textbf{Grand Total Score} & \textbf{80 points (80.0\%)} & \textbf{20 points (20.0\%)} & \textbf{100 points} & \textbf{100\%} \\
\hline
\end{tabularx}

\vspace{1.5mm}
\textbf{Pedagogical Alignment with Course Syllabus (DSAI1003):}
\begin{itemize}[itemsep=0.2pt,topsep=0.2pt]
  \item \textbf{Loops \& Iteration (Part 4):} Definite (\texttt{for}) and indefinite (\texttt{while}) loops, \texttt{range()} arguments, priming reads, sentinel-controlled loops, infinite loop avoidance, \texttt{break}/\texttt{continue}, accumulator and min/max patterns.
  \item \textbf{Functions \& Modular Design (Part 5):} Black-box abstraction, \texttt{def} syntax, parameter passing, return statements vs. \texttt{print()}, default arguments, variable scope (local vs. global shadowing), tuple unpacking, and high cohesion design.
\end{itemize}
\end{tcolorbox}

\vfill
\begin{center}
  \rule{0.5\textwidth}{0.4pt}\\[0.8mm]
  \textbf{\color{primaryblue}--- END OF SOLUTIONS \& MARKING SCHEME ---}
\end{center}

\end{document}
